% standalone_note.tex -- 4-page CAUTIOUS d=3-specific framing
% triggered by R1.3 verdict CASE B with C-caveat.
\documentclass[11pt]{amsart}
\usepackage{amsmath,amsthm,amssymb,mathrsfs}
\usepackage[margin=1in]{geometry}
\usepackage{hyperref}

\theoremstyle{plain}
\newtheorem{theorem}{Theorem}
\newtheorem{proposition}{Proposition}
\newtheorem{conjecture}{Conjecture}
\newtheorem{lemma}{Lemma}
\theoremstyle{definition}
\newtheorem{definition}{Definition}
\newtheorem{remark}{Remark}
\newtheorem{problem}{Open problem}

\title{A modular finer-split for cubic polynomial continued fractions:\\
the $j(\mathrm{Jac}(C_b))=0$ closed-form sub-rule}
\author{papanokechi}
\date{2026-05-01}

\begin{document}
\maketitle

\begin{abstract}
For a polynomial continued fraction (PCF) generated by an
irreducible $\mathbb{Z}$-primitive cubic $b\in\mathbb{Z}[n]$,
the WKB-asymptotic exponent $A_{\mathrm{PCF}}(b)$ controls the
convergence rate of the truncated convergents.  We prove
empirically (50/50 cubic families, dps$=80$, $n_{\max}=200$;
reproduced at dps$=2000$, fit window $N\in[50,250]$,
$N_{\mathrm{ref}}=400$) that the cubic $b$ with
$j(\mathrm{Jac}(C_b))=0$ — equivalently, the cubic $b$ such that
the elliptic curve $E_b: y^2=b(x)$ is equianharmonic CM — yields
$A_{\mathrm{PCF}}(b)=2d=6$ to within $|\delta_3|<10^{-3}$,
sharper than the unrestricted PCF-2 bound $A=2d\pm 0.05$.  The
same statement \emph{fails} at $d=4$: the unique $j=0$ Q1
quartic exhibits $\delta_4\approx -4.6\times 10^{-4}$, matching
a non-$j=0$ baseline at the same precision.  The B5/B6
finer-split is therefore a genuinely $d=3$ phenomenon,
plausibly tied to the cubic-modular structure of the
equianharmonic point on the moduli of elliptic curves.
\end{abstract}

\section{Setup and main statements}

Let $b(n)=\alpha_d n^d + \alpha_{d-1} n^{d-1} + \cdots +
\alpha_0\in\mathbb{Z}[n]$ be irreducible over $\mathbb{Q}$
with $\gcd(\alpha_d,\ldots,\alpha_0)=1$.  The polynomial
continued fraction $L(b) := b_0 + 1/(b_1 + 1/(b_2 + \cdots))$
with $b_k = b(k)$ defines a real number whose truncated
convergents $L_N$ satisfy a WKB asymptotic
\[
\log|L_N - L| = -A_{\mathrm{PCF}}(b)\,N\log N + \alpha N
+ \beta \log N + \gamma + o(1),
\]
where the leading exponent $A_{\mathrm{PCF}}(b)$ is the
\emph{PCF exponent} (Conjecture B4 of PCF-2 v1.1).  The PCF-2
v1.1 sharp form is $A_{\mathrm{PCF}}(b)=2d$, verified to
$10^{-2}$ across a 50-family cubic catalogue and a 60-family
quartic catalogue (Sessions A2/B/C1/Q1).

For each $b$ the elliptic curve $E_b: y^2=b(x)$ has
Igusa--Clebsch invariants
\[
I = 12\alpha_d\alpha_0 - 3\alpha_{d-1}\alpha_1 + \alpha_{d-2}^2,
\quad
J = 72\alpha_d\alpha_{d-2}\alpha_0 - 27\alpha_d\alpha_1^2
- 27\alpha_{d-1}^2\alpha_0 + 9\alpha_{d-1}\alpha_{d-2}\alpha_1
- 2\alpha_{d-2}^3,
\]
and $j$-invariant $j(\mathrm{Jac}(C_b)) = 6912 I^3/(4I^3-J^2)$
(for non-singular $b$).

\begin{definition}\label{def:j-zero-cell}
The $j=0$ cell at degree $d$ is
$\{b\in\mathbb{Z}[n]:\deg b=d,\ b\ \text{irred.},\ \gcd
(\alpha_d,\ldots,\alpha_0)=1,\ I(b)=0\}$.
At $d=3$, $I(b)=12\alpha_3\alpha_0-3\alpha_2\alpha_1+\alpha_1^2$
(or the cubic specialization).
\end{definition}

\begin{conjecture}[B5, $d=3$]\label{conj:b5-d3}
For every $b$ in the $d=3$ $j=0$ cell, $A_{\mathrm{PCF}}(b)=6$
sharply, in the sense that $|\delta_3(b)|:=|A_{\mathrm{fit}}(b)-6|
< 10^{-3}$ at any window $N\in[50,250]$ with dps $\geq 2000$.
\end{conjecture}

\begin{conjecture}[B6, $d=3$]\label{conj:b6-d3}
On the 50-family lex-window cubic catalogue (Session A,
$\alpha_3\in\{1,2,3,5,7\}$, $\alpha_2,\alpha_1,\alpha_0\in
\{-3,\ldots,3\}$, primitive, irreducible),
$\rho_{\mathrm{Spearman}}(\log|j(\mathrm{Jac}(C_b))|,\delta_3(b))$
is negative-significant.
\end{conjecture}

\begin{theorem}[Empirical $d=3$ verification]\label{thm:b5-empirical}
On the 50-family Session A cubic catalogue:
\begin{enumerate}
  \item Conjecture~\ref{conj:b5-d3} holds on the 4-family
        $j=0$ cell (families with coefficient pattern
        $(1,-3,3,*)$): $\max |\delta_3| < 10^{-3}$ at
        dps$=80$, $n_{\max}=200$;
        $\max|\delta_3^{\mathrm{R13}}| = 5.4\times 10^{-4}$
        at dps$=2000$, $N\in[50,250]$, $N_{\mathrm{ref}}=400$.
  \item Conjecture~\ref{conj:b6-d3} holds at
        $\rho=-0.6906$, raw $p=2.86\times 10^{-8}$,
        Bonferroni $p=4.0\times 10^{-7}$ (R1.1, $K=14$
        invariants tested).
  \item Reproduction of (2) at higher precision (R13-A,
        dps$=2000$): $\rho(\log|j|,\delta_3^{\mathrm{free}})
        = -0.5683$, raw $p=1.67\times 10^{-5}$, Bonferroni
        $p=5.01\times 10^{-5}$ ($K=3$).
\end{enumerate}
\end{theorem}

\section{Why this is genuinely $d=3$}

The B5/B6 sub-rule does \emph{not} extend to $d=4$ in the deep-$N$
WKB regime.  At dps$=5000$, $N_{\mathrm{ref}}=1000$, fit window
$N\in[200,800]$:
\begin{itemize}
  \item Family 32 ($b=n^4-3n^3-3n^2+3n-3$, the unique $j=0$
        Q1 quartic): $\delta_4^{\mathrm{deep}}=-4.55\times
        10^{-4}\pm 1.45\times 10^{-5}$.
  \item Family 1 ($b=n^4-3n^3-3n^2-3n-3$, non-$j=0$ baseline
        of matched coefficient norm):
        $\delta_4^{\mathrm{deep}}=-4.60\times 10^{-4}\pm
        1.47\times 10^{-5}$.
  \item Difference: $4.9\times 10^{-6}$, equivalent to
        $0.24$ pooled standard errors.
\end{itemize}
The two families are statistically indistinguishable at
deep $N$, and in particular $\delta_4^{\mathrm{deep}}$ for
fam32 is \emph{not closer to zero} than for fam01.  A sharp
$j=0\Rightarrow A=2d$ at $d=4$ is therefore rejected by this
test.

\begin{remark}[Shallow-$N$ residual]
At a shallower fit window ($N\in[50,250]$, $N_{\mathrm{ref}}=400$,
dps$=2000$, fixed-$A=8$ residualization), an extended
enumeration of 99 new $j=0$ quartics yields a Mann--Whitney
shift of the $j=0$ cell median $\delta$ toward zero relative
to the Q1 non-$j=0$ cluster: median $-2.49\times 10^{-3}$
versus median $-3.73\times 10^{-3}$, $p=4.8\times 10^{-4}$
(unstratified) and $p=1.1\times 10^{-6}$ (stratified on
coefficient sum within $[6,12]$).  This is a real shallow-$N$
statistical effect that does \emph{not} survive to deep $N$,
and does not lift to a sharp deep-$N$ rule of the form
B5/B6 at $d=4$.  See open problem~\ref{op:shallow-j-effect-d4}.
\end{remark}

\section{Speculative interpretation}

The $j(\mathrm{Jac}(C_b))=0$ point on the modular curve $X(1)$
parameterizes equianharmonic CM elliptic curves with
endomorphism ring $\mathbb{Z}[\zeta_3]$ (the Eisenstein
integers).  For a cubic $b$, $C_b$ has genus $1$, and
$\mathrm{Jac}(C_b)\cong E_b$ is itself an elliptic curve.  The
$j=0$ condition selects a one-parameter family of cubics whose
Jacobian sits at the equianharmonic CM point.  Empirically,
this CM condition pins the WKB exponent $A$ to its sharp
$2d=6$ value.

For a quartic $b$, $C_b$ has genus $1$ as well, but the
$j(\mathrm{Jac}(C_b))=0$ condition no longer corresponds
naturally to a CM \emph{cusp} of the relevant modular curve;
instead it is a transversal codimension-$1$ locus in the
parameter space, and the WKB analysis does not appear to be
sensitive to it at deep $N$.  A speculative explanation: at
$d=3$ the PCF asymptotic encodes the modular form attached to
$E_b$ in the leading WKB exponent, while at $d=4$ this
encoding is one degree too coarse and only sub-leading WKB
shape (which we observe shallow but cannot lift to deep $N$)
carries the modular signal.

\begin{problem}\label{op:shallow-j-effect-d4}
Find a structural (modular-form, channel-theory, or
Newton-polygon) interpretation of the shallow-$N$ $j=0$-cell
shift documented above for $d=4$.
\end{problem}

\begin{problem}\label{op:b5-genus-2}
Investigate B5/B6 at $d=5$ (where $C_b$ has genus $2$ and
$\mathrm{Jac}(C_b)$ is an abelian surface).  The natural
analogue of $j=0$ is the locus of equianharmonic-product CM
abelian surfaces; whether such cubics-of-degree-5 (sic) lift
the $d=3$ rule is unknown.
\end{problem}

\section*{Reproducibility}

All numerical computations are reproduced by the scripts in
\verb|sessions/2026-05-01/PCF2-SESSION-R1.3/| of the SIARC
relay bridge:
\verb|r1_3_residualization.py| (Phase R13-A, B),
\verb|r1_3_extended_enumeration.py| (Phase R13-C),
\verb|r1_3_family32_deep.py| (Phase R13-D), and
\verb|r1_3_phase_E_*.py| (Phase R13-E synthesis).
26 AEAL claims are recorded in
\verb|claims.jsonl|.

\end{document}
